\honest{Epistemic Viciousness}{Eliezer Yudkowsky}{2009}

Someone pointed me to a twelve-page essay by Gillian Russell, ``Epistemic Viciousness in the Martial Arts''; I cannot remember who. Russell tells of believing, at eleven, a teacher's promise that three years of practising punches in the air would make it possible to ``kill a bull with one blow''; she expected to manage it by fourteen. By ``viciousness'' the essay means ``possessing of vices''.

``It all generalizes amazingly.'' I summarize her account. The art, the dojo and the teacher are sacred, and other styles are ``the wrong religion''. Someone whose teacher took them aside to teach a special move, and who has practised it for twenty years, will want to discard evidence against it. A beginner cannot learn from a book and has to trust the teacher. Historical masters are deferred to, whereas runners do not defer to the ancient Greeks, and a physicist may say that Newton's theories are false. Martial artists have too few data to test their techniques, and so cannot test their training either: should you practise a strike in the air, or will that teach you to overextend? Worst, they ``live in poverty but continue to act as though we live in luxury''. \nb{The summary follows Russell's sections in order, and its quotations match the draft on Russell's website closely.}

I remembered one more point from the essay, but on a second reading it was not there: martial arts declined when real fights stopped, and now there are videos of black belts ``being pounded into the ground by someone with real fighting experience.'' \nb{The essay makes a nearby point: judo leaves ``relatively little room for pretence'', since a throw either works or does not, and odd beliefs do better in arts with less competitive sparring. For the history of decline, the post gives no source, and videos of defeats today do not show that things were once better.} I also recall, without the link, a ki master whose students fell at his touchless attacks, and who was beaten by a skeptic. \nb{Probably Yanagi Ryuken, who offered money to any mixed martial artist who could beat him, and lost on Japanese television. A commenter linked the video two days later.}

``How to not lose'' is more widely useful than ``how to win.'' ``Every single one of these risk factors transfers straight over to any attempt to start a `rationality dojo'.'' I do not go through them one by one. I ask: ``What can be done about it?'' \nb{Russell's essay has partial answers: test parts of a technique when the whole cannot be tested, and keep justified deference to a teacher apart from the unjustified kind.}
